% Plantilla del informe de la Parte 1 (Unidad 1) - Proyecto integrador
% Compilar con: pdflatex -> biber -> pdflatex -> pdflatex (o latexmk -pdf)
\documentclass[12pt,a4paper]{article}
\usepackage[utf8]{inputenc}
\usepackage[T1]{fontenc}
\usepackage[spanish,es-nodecimaldot]{babel}
\usepackage{times}
\usepackage[margin=2.54cm]{geometry}
\usepackage{setspace}
\doublespacing
\usepackage{csquotes}
\usepackage[style=apa,backend=biber]{biblatex}
\DeclareLanguageMapping{spanish}{spanish-apa}
\addbibresource{referencias.bib}
\usepackage{booktabs,longtable,array}
\usepackage{graphicx}
\usepackage{caption}
\usepackage{hyperref}
\usepackage{parskip}
\setlength{\parindent}{0.5in}
\setlength{\parskip}{0pt}

\begin{document}

% ---------------- PORTADA ----------------
\begin{titlepage}
\centering
\singlespacing
\vspace*{2cm}
{\bfseries T\'itulo del proyecto (m\'aximo 12 palabras)\par}
\vspace{1.5cm}
Equipo N.$^{\circ}$ \_\_\ \\
Integrantes: Nombre Apellido, Nombre Apellido, Nombre Apellido, Nombre Apellido\\[1cm]
Carrera de Ingenier\'ia Ambiental, Universidad Nacional de Loja\\
Programaci\'on y Base de Datos\\
Docente: Carlos Guillermo Chuncho Morocho\\[1cm]
Octubre de 2026
\end{titlepage}

% ---------------- RESUMEN ----------------
\section*{Resumen}
\singlespacing
Resumen de 120 a 150 palabras: tema, zona, problema, pregunta, hip\'otesis y objetivo general.

\noindent\textit{Palabras clave:} FIRELAB-Loja, tema, zona, programaci\'on, Python.
\doublespacing
\newpage

% ---------------- 1. INTRODUCCION ----------------
\section{Introducci\'on y contexto}
Contexto territorial del cant\'on Loja, descripci\'on breve del proyecto FIRELAB-Loja y relevancia del tema del equipo. Incluya citas verificables \parencite{firelab2025}.

% ---------------- 2. PROBLEMA ----------------
\section{Planteamiento del problema}
Situaci\'on, vac\'io o dificultad y oportunidad de resolverlo con un programa y los datos asignados.

% ---------------- 3. PREGUNTA ----------------
\section{Pregunta de investigaci\'on}
Una sola pregunta: qu\'e se analiza, d\'onde (zona) y cu\'ando (2019--2024).

% ---------------- 4. HIPOTESIS ----------------
\section{Hip\'otesis}
Afirmaci\'on contrastable, con variables, zona, periodo y criterio cuantitativo.

% ---------------- 5. OBJETIVOS ----------------
\section{Objetivos}
\subsection{Objetivo general}
Verbo en infinitivo + qu\'e + d\'onde + con qu\'e datos.

\subsection{Objetivos espec\'ificos}
\begin{enumerate}
  \item Objetivo espec\'ifico 1.
  \item Objetivo espec\'ifico 2.
  \item Objetivo espec\'ifico 3 (opcional; m\'aximo tres).
\end{enumerate}

% ---------------- 6. DATOS ----------------
\section{Datos y variables}
Base asignada, periodo (2019--2024), zona y unidad de observaci\'on (una celda de 500\,m en un mes). Variables:

\begin{table}[ht]
\centering
\singlespacing
\caption*{\textbf{Tabla 1}\\\textit{Variables seleccionadas para el proyecto}}
\begin{tabular}{llll}
\toprule
Variable & Significado & Unidad & Tipo \\
\midrule
nombre\_1 & ... & ... & num\'erica \\
nombre\_2 & ... & ... & categ\'orica \\
\bottomrule
\end{tabular}
\end{table}
\doublespacing

% ---------------- 7. ALCANCE ----------------
\section{Alcance y limitaciones}
Qu\'e incluye y qu\'e no (por ejemplo, el a\~no 2025 no se utiliza).

% ---------------- 8. MATRIZ ----------------
\section{Matriz de coherencia}
\begin{table}[ht]
\centering
\singlespacing
\caption*{\textbf{Tabla 2}\\\textit{Matriz de coherencia de la Parte 1}}
\begin{tabular}{p{3cm}p{7cm}p{4cm}}
\toprule
Elemento & Redacci\'on & Salida verificable \\
\midrule
Problema & & \\
Pregunta & & \\
Hip\'otesis & & \\
Objetivo general & & \\
OE1 & & \\
OE2 & & \\
OE3 & & \\
\bottomrule
\end{tabular}
\end{table}
\doublespacing

% ---------------- REFERENCIAS ----------------
\printbibliography[title={Referencias}]

% ---------------- ANEXO ----------------
\appendix
\section{Uso de inteligencia artificial}
Declare para qu\'e se us\'o, qu\'e se verific\'o y qu\'e parte fue reescrita por el equipo.

\end{document}
